% -- Radiopaedia Brain CT Scraper: Source Code Listings --
% -- Gunakan minipage agar listing tidak terpotong ke halaman berikutnya --



\begin{minipage}{\textwidth}
\begin{lstlisting}[caption={Import library dan konstanta global (scraper\_base.py)}, label={lst:base_imports}, frame=single, captionpos=t, language=Python]
import io, os, re, sys, json, time, random, logging, hashlib, shutil, requests
from pathlib import Path
from urllib.parse import urlencode, urljoin, urlparse

from bs4 import BeautifulSoup
from tqdm import tqdm
from PIL import Image
from io import BytesIO
from selenium import webdriver
from selenium.webdriver.common.by import By
from selenium.webdriver.common.keys import Keys
from selenium.webdriver.chrome.service import Service
from selenium.webdriver.chrome.options import Options
from selenium.webdriver.support.ui import WebDriverWait
from selenium.webdriver.support import expected_conditions as EC
from selenium.common.exceptions import (
    TimeoutException, NoSuchElementException,
    WebDriverException, StaleElementReferenceException,
    ElementClickInterceptedException,
)
from webdriver_manager.chrome import ChromeDriverManager

import config

COOKIE_SYNC_INTERVAL = 900
logger = logging.getLogger(__name__)
\end{lstlisting}
\end{minipage}


\begin{minipage}{\textwidth}
\begin{lstlisting}[caption={Konfigurasi logging ke file dan stdout (scraper\_base.py)}, label={lst:setup_logging}, frame=single, captionpos=t, language=Python]
def setup_logging(log_file: str) -> None:
    try:
        stream = io.TextIOWrapper(
            sys.stdout.buffer, encoding="utf-8",
            errors="replace", line_buffering=True
        )
    except AttributeError:
        stream = sys.stdout
    logging.basicConfig(
        level=logging.DEBUG,
        format="%(asctime)s [%(levelname)s]  %(message)s",
        handlers=[
            logging.FileHandler(log_file, encoding="utf-8"),
            logging.StreamHandler(stream),
        ],
    )
\end{lstlisting}
\end{minipage}


\begin{minipage}{\textwidth}
\begin{lstlisting}[caption={Fungsi-fungsi utilitas pembantu (scraper\_base.py)}, label={lst:utils}, frame=single, captionpos=t, language=Python]
def random_delay(min_s=None, max_s=None):
    delay = random.uniform(
        min_s if min_s is not None else config.REQUEST_DELAY_MIN,
        max_s if max_s is not None else config.REQUEST_DELAY_MAX,
    )
    logger.debug(f"Sleeping {delay:.2f}s ...")
    time.sleep(delay)


def make_output_dir(category: str) -> Path:
    p = Path(config.OUTPUT_DIR) / category
    p.mkdir(parents=True, exist_ok=True)
    logger.info(f"Output directory ready: '{p}'")
    return p


def image_hash(data: bytes) -> str:
    return hashlib.md5(data).hexdigest()


def is_ct_modality(text: str) -> bool:
    return any(mod in text.lower() for mod in config.ALLOWED_MODALITIES)


def sanitize_filename(name: str) -> str:
    return re.sub(r'[\\/*?:"<>|]', "_", name)


def is_driver_alive(driver: webdriver.Chrome) -> bool:
    try:
        _ = driver.current_url
        return True
    except Exception:
        return False
\end{lstlisting}
\end{minipage}


\begin{minipage}{\textwidth}
\begin{lstlisting}[caption={Inisialisasi ChromeDriver dengan opsi anti-deteksi (scraper\_base.py)}, label={lst:build_driver}, frame=single, captionpos=t, language=Python]
DOWNLOAD_DIR = Path(config.OUTPUT_DIR) / "_downloads"

def build_driver() -> webdriver.Chrome:
    DOWNLOAD_DIR.mkdir(parents=True, exist_ok=True)

    opts = Options()
    opts.page_load_strategy = "eager"
    if config.HEADLESS_BROWSER:
        opts.add_argument("--headless=new")
    opts.add_argument("--window-size=1920,1080")
    opts.add_argument("--disable-blink-features=AutomationControlled")
    opts.add_experimental_option("excludeSwitches", ["enable-automation"])
    opts.add_experimental_option("useAutomationExtension", False)
    opts.add_argument(
        "user-agent=Mozilla/5.0 (Windows NT 10.0; Win64; x64) "
        "AppleWebKit/537.36 (KHTML, like Gecko) Chrome/122.0.0.0 Safari/537.36"
    )
    prefs = {
        "download.default_directory": str(DOWNLOAD_DIR.resolve()),
        "download.prompt_for_download": False,
        "download.directory_upgrade": True,
        "safebrowsing.enabled": True,
    }
    opts.add_experimental_option("prefs", prefs)

    driver = webdriver.Chrome(
        service=Service(ChromeDriverManager().install()), options=opts
    )
    driver.set_page_load_timeout(config.PAGE_LOAD_TIMEOUT)
    driver.execute_script(
        "Object.defineProperty(navigator, 'webdriver', {get: () => undefined})"
    )
    return driver
\end{lstlisting}
\end{minipage}


\begin{minipage}{\textwidth}
\begin{lstlisting}[caption={Proses login ke Radiopaedia dengan multi-selector fallback (scraper\_base.py)}, label={lst:login}, frame=single, captionpos=t, language=Python]
def login(driver: webdriver.Chrome) -> bool:
    logger.info("Navigating to login page ...")
    driver.get("https://radiopaedia.org/sessions/new")
    try:
        WebDriverWait(driver, config.PAGE_LOAD_TIMEOUT).until(
            lambda d: d.execute_script("return document.readyState") == "complete"
        )
    except TimeoutException:
        pass
    random_delay(4, 6)
    _dismiss_cookie_banner(driver)

    EMAIL_SELECTORS = [
        (By.NAME, "user[identity]"), (By.ID, "user_identity"),
        (By.NAME, "session[email]"), (By.CSS_SELECTOR, "input[type='email']"),
        (By.XPATH, "//input[@autocomplete='email']"),
    ]
    email_field = _find_first(driver, EMAIL_SELECTORS, timeout=15)
    if email_field is None:
        email_field = _find_visible_input(driver, "email", timeout=10)
    if email_field is None:
        logger.error("Login FAILED: email input field tidak ditemukan.")
        return False
    email_field.clear()
    email_field.send_keys(config.RADIOPAEDIA_EMAIL)
    random_delay(0.5, 1.2)

    PASSWORD_SELECTORS = [
        (By.NAME, "user[password]"), (By.ID, "user_password"),
        (By.CSS_SELECTOR, "input[type='password']"),
        (By.XPATH, "//input[@autocomplete='current-password']"),
    ]
    password_field = _find_first(driver, PASSWORD_SELECTORS, timeout=10)
    if password_field is None:
        password_field = _find_visible_input(driver, "password", timeout=10)
    if password_field is None:
        logger.error("Login FAILED: password input field tidak ditemukan.")
        return False
    password_field.clear()
    password_field.send_keys(config.RADIOPAEDIA_PASSWORD)
    random_delay(0.5, 1.2)

    submit_btn = _find_first(driver, [
        (By.CSS_SELECTOR, "input[type='submit']"),
        (By.XPATH, "//button[contains(text(),'Sign in')]"),
        (By.XPATH, "//button[contains(text(),'Log in')]"),
    ], timeout=10)
    if submit_btn:
        submit_btn.click()
    else:
        password_field.send_keys(Keys.RETURN)
    random_delay(4, 7)

    logged_in = _find_first(driver, [
        (By.CSS_SELECTOR, "a[href*='sign_out']"),
        (By.XPATH, "//*[contains(@class,'user-menu')]"),
        (By.XPATH, "//a[contains(text(),'Sign out')]"),
    ], timeout=12)
    if logged_in:
        logger.info("[OK] Login successful.")
        return True

    logger.error("Login FAILED: tidak terdeteksi status logged-in.")
    return False
\end{lstlisting}
\end{minipage}


\begin{minipage}{\textwidth}
\begin{lstlisting}[caption={Pembangun URL pencarian dan pembaca jumlah halaman (scraper\_base.py)}, label={lst:search_pagination}, frame=single, captionpos=t, language=Python]
def build_search_url(query: str, page: int = 1) -> str:
    params = [
        ("scope", "cases"),
        ("q", query),
        ("page", page),
        ("modality[]", "CT"),
    ]
    return f"https://radiopaedia.org/search?{urlencode(params)}"


def get_total_pages(driver: webdriver.Chrome) -> int:
    try:
        soup = BeautifulSoup(driver.page_source, "lxml")
        for selector in [
            "ul.pagination li a", "nav.pagination a",
            "div.pagination a", "[class*='pagina'] a"
        ]:
            nums = [
                int(a.get_text(strip=True))
                for a in soup.select(selector)
                if a.get_text(strip=True).isdigit()
            ]
            if nums:
                return max(nums)
        if soup.select_one("a[rel='next']"):
            return 99
        return 1
    except Exception as e:
        logger.warning(f"Could not determine total pages: {e}")
        return 1


def _is_relevant_case(title: str, keyword: str) -> bool:
    STOPWORDS = {"and","or","the","of","a","an","in","with","for","to","on","ct","mri"}
    kw_words = [
        w for w in re.split(r'\W+', keyword.lower())
        if w and w not in STOPWORDS
    ]
    if not kw_words:
        return True
    return any(w in title.lower() for w in kw_words)
\end{lstlisting}
\end{minipage}


\begin{minipage}{\textwidth}
\begin{lstlisting}[caption={Ekstraksi tautan kasus dari halaman hasil pencarian (scraper\_base.py)}, label={lst:get_case_links}, frame=single, captionpos=t, language=Python]
def get_case_links_from_page(driver: webdriver.Chrome, keyword: str = "") -> list:
    soup = BeautifulSoup(driver.page_source, "lxml")
    BASE = "https://radiopaedia.org"
    links = []

    result_cards = soup.select(
        "a.search-result.case, a.search-result[href*='/cases/'], "
        "div.search-result a[href*='/cases/'], "
        "article.search-result a[href*='/cases/']"
    )

    if result_cards:
        for a in result_cards:
            href = a.get("href", "")
            if not href or "/cases/" not in href:
                continue
            parent  = a.find_parent(["li", "div", "article"])
            heading = a.find(["h2", "h3", "h4", "strong"])
            if not heading and parent:
                heading = parent.find(["h2", "h3", "h4", "strong"])
            title = heading.get_text(strip=True) if heading else a.get_text(strip=True)
            if keyword and title and not _is_relevant_case(title, keyword):
                logger.debug(f"  [filter] Dilewati: '{title}'")
                continue
            url = urljoin(BASE, href)
            if url not in links:
                links.append(url)

    if not links:
        for selector in [
            "div[class*='search'] a[href*='/cases/']",
            "div[class*='result'] a[href*='/cases/']",
            "a[href*='/cases/']",
        ]:
            for a in soup.select(selector):
                url = urljoin(BASE, a["href"])
                if url not in links:
                    links.append(url)
            if links:
                break

    logger.debug(f"  Found {len(links)} case link(s).")
    return links
\end{lstlisting}
\end{minipage}


\begin{minipage}{\textwidth}
\begin{lstlisting}[caption={Sinkronisasi cookie dan restart driver (scraper\_base.py)}, label={lst:cookie_sync}, frame=single, captionpos=t, language=Python]
def sync_cookies(driver, session) -> bool:
    try:
        session.cookies.clear()
        for cookie in driver.get_cookies():
            session.cookies.set(cookie["name"], cookie["value"])
        logger.debug("Cookies synced from browser to requests session.")
        return True
    except Exception as e:
        logger.warning(f"sync_cookies failed: {e}")
        return False


def restart_driver_and_login(seen_case_urls, seen_hashes) -> tuple:
    logger.warning("Restarting browser session ...")
    try:
        driver = build_driver()
    except Exception as e:
        raise RuntimeError(f"Could not rebuild driver: {e}") from e
    session = requests.Session()
    if not login(driver):
        driver.quit()
        raise RuntimeError("Re-login failed after browser restart.")
    sync_cookies(driver, session)
    logger.info("Browser restarted and logged in successfully.")
    return driver, session
\end{lstlisting}
\end{minipage}


\begin{minipage}{\textwidth}
\begin{lstlisting}[caption={Orkestrator scraping per kategori (scraper\_base.py)}, label={lst:scrape_category}, frame=single, captionpos=t, language=Python]
def scrape_category(
    driver, session, category, keywords, output_dir,
    target_count, seen_hashes, seen_case_urls,
) -> tuple:
    downloaded = sum(1 for _ in output_dir.iterdir())
    logger.info(
        f"\n{'='*60}\nCategory : {category.upper()}\n"
        f"Target   : {target_count} images | Already have: {downloaded}\n{'='*60}"
    )
    pbar = tqdm(total=target_count, initial=downloaded, desc=category, unit="img")
    last_cookie_sync = time.time()

    def _ensure_driver_alive():
        nonlocal driver, session, last_cookie_sync
        if not is_driver_alive(driver):
            logger.error("Browser session is dead -- attempting restart.")
            try:
                driver.quit()
            except Exception:
                pass
            driver, session = restart_driver_and_login(seen_case_urls, seen_hashes)
            last_cookie_sync = time.time()

    def _maybe_sync_cookies():
        nonlocal last_cookie_sync
        now = time.time()
        if now - last_cookie_sync >= COOKIE_SYNC_INTERVAL:
            _ensure_driver_alive()
            if sync_cookies(driver, session):
                last_cookie_sync = now

    for keyword in keywords:
        if downloaded >= target_count:
            break
        logger.info(f"  Searching: '{keyword}'")
        page = 1
        _ensure_driver_alive()
        driver.get(build_search_url(keyword, page=1))
        random_delay()
        total_pages = get_total_pages(driver)

        while page <= total_pages and downloaded < target_count:
            driver.get(build_search_url(keyword, page=page))
            random_delay()
            case_links = get_case_links_from_page(driver, keyword=keyword)
            if not case_links:
                page += 1
                continue

            for case_url in case_links:
                if downloaded >= target_count:
                    break
                if case_url in seen_case_urls:
                    continue
                seen_case_urls.add(case_url)
                _maybe_sync_cookies()

                try:
                    n = download_images_from_case(
                        driver=driver, session=session,
                        case_url=case_url, output_dir=output_dir,
                        seen_hashes=seen_hashes,
                        target_count=target_count,
                        current_count=downloaded,
                    )
                except Exception as e:
                    logger.error(f"  Error di case {case_url}: {e}")
                    _ensure_driver_alive()
                    continue

                if n > 0:
                    downloaded += n
                    pbar.update(n)
                    logger.info(
                        f"  [RESULT] {case_url} -> {n} gambar. "
                        f"Total: {downloaded}/{target_count}"
                    )
                random_delay()
            page += 1

    pbar.close()
    logger.info(f"  Finished '{category}': {downloaded} images collected.")
    return downloaded, driver, session
\end{lstlisting}
\end{minipage}


\begin{minipage}{\textwidth}
\begin{lstlisting}[caption={Scraper kategori glioma (scrape\_glioma.py)}, label={lst:scrape_glioma}, frame=single, captionpos=t, language=Python]
"""
scrape_glioma.py
----------------
Scraper untuk kategori GLIOMA dari Radiopaedia.
Jalankan langsung:  python scrape_glioma.py
"""

import requests
from scraper_base import make_output_dir, build_driver, login, sync_cookies, scrape_category
import config

CATEGORY = "glioma"
TARGET   = config.TARGET_PER_CATEGORY.get(CATEGORY, 300)
KEYWORDS = config.CATEGORIES.get(CATEGORY, [
    "glioma brain CT",
    "glioblastoma CT head",
    "glioblastoma multiforme CT",
    "astrocytoma brain CT",
    "high grade glioma CT",
    "low grade glioma CT",
    "diffuse glioma CT",
    "anaplastic astrocytoma CT",
    "oligodendroglioma CT",
])


def main():
    output_dir = make_output_dir(CATEGORY)
    driver     = build_driver()
    session    = requests.Session()

    try:
        if not login(driver):
            print("Login gagal. Periksa kredensial di config.py.")
            return

        sync_cookies(driver, session)
        count, driver, session = scrape_category(
            driver=driver,
            session=session,
            category=CATEGORY,
            keywords=KEYWORDS,
            output_dir=output_dir,
            target_count=TARGET,
            seen_hashes=set(),
            seen_case_urls=set(),
        )
        print(f"Selesai -- {CATEGORY.upper()}: {count} gambar terkumpul.")

    except KeyboardInterrupt:
        print("\nDihentikan oleh pengguna.")
    except Exception as e:
        print(f"Error tak terduga: {e}")
    finally:
        try:
            driver.quit()
        except Exception:
            pass


if __name__ == "__main__":
    main()
\end{lstlisting}
\end{minipage}


\begin{minipage}{\textwidth}
\begin{lstlisting}[caption={Scraper kategori meningioma (scrape\_meningioma.py)}, label={lst:scrape_meningioma}, frame=single, captionpos=t, language=Python]
"""
scrape_meningioma.py
--------------------
Scraper untuk kategori MENINGIOMA dari Radiopaedia.
Jalankan langsung:  python scrape_meningioma.py
"""

import requests
from scraper_base import make_output_dir, build_driver, login, sync_cookies, scrape_category
import config

CATEGORY = "meningioma"
TARGET   = config.TARGET_PER_CATEGORY.get(CATEGORY, 300)
KEYWORDS = config.CATEGORIES.get(CATEGORY, [
    "meningioma CT head",
    "meningioma brain CT",
    "meningioma CT scan",
    "convexity meningioma CT",
    "sphenoid wing meningioma CT",
    "parasagittal meningioma CT",
    "falcine meningioma CT",
    "skull base meningioma CT",
    "olfactory groove meningioma CT",
    "posterior fossa meningioma CT",
])


def main():
    output_dir = make_output_dir(CATEGORY)
    driver     = build_driver()
    session    = requests.Session()

    try:
        if not login(driver):
            print("Login gagal. Periksa kredensial di config.py.")
            return

        sync_cookies(driver, session)
        count, driver, session = scrape_category(
            driver=driver,
            session=session,
            category=CATEGORY,
            keywords=KEYWORDS,
            output_dir=output_dir,
            target_count=TARGET,
            seen_hashes=set(),
            seen_case_urls=set(),
        )
        print(f"Selesai -- {CATEGORY.upper()}: {count} gambar terkumpul.")

    except KeyboardInterrupt:
        print("\nDihentikan oleh pengguna.")
    except Exception as e:
        print(f"Error tak terduga: {e}")
    finally:
        try:
            driver.quit()
        except Exception:
            pass


if __name__ == "__main__":
    main()
\end{lstlisting}
\end{minipage}


\begin{minipage}{\textwidth}
\begin{lstlisting}[caption={Scraper kategori normal (scrape\_normal.py)}, label={lst:scrape_normal}, frame=single, captionpos=t, language=Python]
"""
scrape_normal.py
----------------
Scraper untuk kategori NORMAL (no tumor) dari Radiopaedia.
Jalankan langsung:  python scrape_normal.py
"""

import requests
from scraper_base import make_output_dir, build_driver, login, sync_cookies, scrape_category
import config

CATEGORY = "normal"
TARGET   = config.TARGET_PER_CATEGORY.get(CATEGORY, 300)
KEYWORDS = config.CATEGORIES.get(CATEGORY, [
    "normal head CT",
    "normal brain CT",
    "normal CT head adult",
    "normal CT head pediatric",
    "normal CT head neonate",
    "normal CT brain anatomy",
    "normal CT head and cervical spine",
    "normal CTA head",
    "normal head CT venogram",
    "normal head CT perfusion",
])


def main():
    output_dir = make_output_dir(CATEGORY)
    driver     = build_driver()
    session    = requests.Session()

    try:
        if not login(driver):
            print("Login gagal. Periksa kredensial di config.py.")
            return

        sync_cookies(driver, session)
        count, driver, session = scrape_category(
            driver=driver,
            session=session,
            category=CATEGORY,
            keywords=KEYWORDS,
            output_dir=output_dir,
            target_count=TARGET,
            seen_hashes=set(),
            seen_case_urls=set(),
        )
        print(f"Selesai -- {CATEGORY.upper()}: {count} gambar terkumpul.")

    except KeyboardInterrupt:
        print("\nDihentikan oleh pengguna.")
    except Exception as e:
        print(f"Error tak terduga: {e}")
    finally:
        try:
            driver.quit()
        except Exception:
            pass


if __name__ == "__main__":
    main()
\end{lstlisting}
\end{minipage}


\begin{minipage}{\textwidth}
\begin{lstlisting}[caption={Scraper kategori pituitary (scrape\_pituitary.py)}, label={lst:scrape_pituitary}, frame=single, captionpos=t, language=Python]
"""
scrape_pituitary.py
-------------------
Scraper untuk kategori PITUITARY (tumor hipofisis) dari Radiopaedia.
Jalankan langsung:  python scrape_pituitary.py
"""

import requests
from scraper_base import make_output_dir, build_driver, login, sync_cookies, scrape_category
import config

CATEGORY = "pituitary"
TARGET   = config.TARGET_PER_CATEGORY.get(CATEGORY, 300)
KEYWORDS = config.CATEGORIES.get(CATEGORY, [
    "pituitary adenoma CT",
    "pituitary tumor CT",
    "pituitary macroadenoma CT",
    "pituitary microadenoma CT",
    "sella turcica tumor CT",
    "craniopharyngioma CT",
    "pituitary apoplexy CT",
    "sellar mass CT",
    "suprasellar mass CT",
    "pituitary CT head",
])


def main():
    output_dir = make_output_dir(CATEGORY)
    driver     = build_driver()
    session    = requests.Session()

    try:
        if not login(driver):
            print("Login gagal. Periksa kredensial di config.py.")
            return

        sync_cookies(driver, session)
        count, driver, session = scrape_category(
            driver=driver,
            session=session,
            category=CATEGORY,
            keywords=KEYWORDS,
            output_dir=output_dir,
            target_count=TARGET,
            seen_hashes=set(),
            seen_case_urls=set(),
        )
        print(f"Selesai -- {CATEGORY.upper()}: {count} gambar terkumpul.")

    except KeyboardInterrupt:
        print("\nDihentikan oleh pengguna.")
    except Exception as e:
        print(f"Error tak terduga: {e}")
    finally:
        try:
            driver.quit()
        except Exception:
            pass


if __name__ == "__main__":
    main()
\end{lstlisting}
\end{minipage}
